\subsection{页面识别与证据记录}

处理流程以论文为单位建立 JSONL 页面记录。唯一可用文本层的论文直接提取，其余
58 篇先按 120 DPI 渲染并由 RapidOCR 全量识别；文本密集页出现 OCR 中位置信度
低于 0.75、中文正文少于 100 字或标题段落明显断裂时，按 220 DPI 重渲染，并以
Tesseract 的 \texttt{chi\_sim+eng} 复识。每页记录年份、题型、论文编号、页码、
识别方法、置信度、中文字符数、文本框坐标及复识审计信息。

记录表把证据细分为章节命中、模型链、写作特征和页面视觉核验四类。公式、表格及
小字号图注不依据 OCR 猜测；这类内容只通过原始栅格的局部裁剪核对。统计脚本必须
先验证 59 个文件和 2892 个连续页码，再计算章节顺序、页数分布、段落密度、图表
候选、公式候选和模型词链。自动页数统计得到全篇页数中位数为 45 页，四分位范围为
37--58.5 页；每篇论文的页面文本行数中位数再取中位数为 38，段落候选数对应
中位数为 6。段落候选
由正文行缩进和垂直间距估计，只用于版面密度比较，不等同于排版软件中的自然段。
共有 2232 页触发并完成 Tesseract 复识，其中 1583 页因质量分更高而被选为最终
文本；最终方法还包括 1281 页 RapidOCR 和 28 页直接文本提取。按同一严格阈值，
1770 页在严格的逐字 OCR 质量阈值下仍标记为未确认；这表示其机器文本不用于逐字引文，
不表示页面未读。59 篇均已逐页人工通读，涉及公式、表格和小字号图注的结论另由原始栅格
或局部裁剪确认。
十四类章节的检出率和页级跨度见表~\ref{tab:section-stats}。

